\documentclass[11pt]{article}
\usepackage[margin=1.1in]{geometry}
\usepackage{amsmath,amssymb,amsthm}
\usepackage{booktabs}
\usepackage[colorlinks=true,linkcolor=blue,citecolor=blue,urlcolor=blue]{hyperref}

\newtheorem{theorem}{Theorem}[section]
\newtheorem{lemma}[theorem]{Lemma}
\newtheorem{proposition}[theorem]{Proposition}
\newtheorem{corollary}[theorem]{Corollary}
\theoremstyle{remark}
\newtheorem{remark}[theorem]{Remark}

\newcommand{\F}{\mathbb{F}}
\newcommand{\Dfour}{\mathsf{D4}}
\newcommand{\chis}{\chi_\sigma}
\newcommand{\chip}{\chi_P}

\title{The Cohomological Price of NOT}
\author{Won Chul Yang\\
\small independent researcher \quad \texttt{wcy0969@gmail.com}}
\date{2026}

\begin{document}
\maketitle

\begin{abstract}
The minimal number of negations needed to compute a Boolean function is
governed by Markov's decrease invariant $d(f)$: circuits need
$\lceil\log_2(d(f)+1)\rceil$ negations (Markov 1958), and formulas need
exactly $d(f)$ (Morizumi 2009). We show that the decrease invariant carries a
\emph{cohomological address}. On the dual-rail four-valued lattice, the
negation $\sigma$ and the rail-swap $P$ generate a Klein four-group $V_4$
whose character group $\operatorname{Hom}(V_4,\F_2)\cong\F_2^2$ has exactly
three nonzero characters, and $\chis$ is the \emph{unique} one that prices
negation --- the two alternatives fail in opposite directions, witnessed by
$P$ (face-flip $1$, cost $0$) and $\sigma P$ (face-flip $0$, cost $1$). The
main theorem reads the full invariant as an iterated extension: $d(f)$ equals
the minimal length of a nested $\chis$-tower, each rung adjoining exactly one
negation inside a monotone context. The upper bound uses a single uniform
three-zone witness; the lower bound is an alternation argument valid over any
finite poset. All statements are machine-verified in core Lean~4 with kernel
axioms \texttt{propext} and \texttt{Quot.sound} only.
\end{abstract}

\section{Introduction}

Markov \cite{markov} determined the inversion complexity of Boolean
circuits: $I(f)=\lceil\log_2(d(f)+1)\rceil$, where the decrease $d(f)$ counts
the maximal number of $1\to 0$ descents of $f$ along an increasing chain of
inputs. Morizumi \cite{morizumi} settled the formula model:
$I_{\mathrm{for}}(f)=d(f)$ exactly. For multi-valued logic,
Kochergin--Mikhailovich \cite{km} extended the \emph{circuit} law (still
logarithmic); no formula-exact multi-valued law appears in the literature.

This note contributes three things.

\textbf{(i) A four-valued frame.} On the dual-rail lattice $\Dfour$ the
exact formula-model law persists (the companion preprint \cite{wcy2} proves
$\nu(f)=d(f)$ in that setting; we do not reprove it here).

\textbf{(ii) A cohomological reading.} The negation $\sigma$ and the
monotone rail-swap $P$ generate $V_4$, and among the three nonzero
characters of $H^1(V_4;\F_2)$ exactly one --- $\chis$ --- prices negation
(\S\ref{sec:char}). The seemingly natural candidate, the face-flip character
$\chis+\chip$, fails in \emph{both} directions.

\textbf{(iii) A nested normal form.} $d(f)$ is the minimal length of a
nested tower $T_0\subseteq T_1\subseteq\cdots$ in which each rung adjoins
exactly one negation inside a monotone context (\S\ref{sec:tower}).
Morizumi's upper-bound construction is the \emph{parallel} decomposition
$f_1\vee(f_2\wedge\neg f_t)$; the nested single-peel form proved here is
strictly more structured, and appears to be new even in the Boolean case.
The upper bound is produced by one uniform three-zone witness
(\S\ref{sec:witness}), whose correctness is a short budget argument; a
cautionary example shows the obvious witness is wrong already at $d(f)=1$.

Everything below is machine-verified in core Lean~4 (\S\ref{sec:lean}).

\section{Preliminaries}

Let $R=\{0,1\}$ with $0<1$, and order $R^n$ componentwise. A \emph{chain} is
an increasing sequence $x^0<x^1<\cdots<x^m$ in $R^n$. For $f:R^n\to R$ and a
chain $X$, the \emph{decrease} $d_X(f)$ is the number of indices $i$ with
$f(x^i)=1$ and $f(x^{i+1})=0$; $d(f)=\max_X d_X(f)$ over maximal chains. For
a binary word $w$, write $\mathrm{rises}(w)$ and $\mathrm{drops}(w)$ for the
numbers of $01$- and $10$-adjacencies. The \emph{decrease potential} is
\[
\omega(x)\;:=\;\max\{\,\text{drops of }f\text{ along an increasing chain
ending at }x\,\},
\]
so $\omega$ is nondecreasing along chains, $\omega\le d(f)$ everywhere, and
every descent step strictly increases $\omega$.

The dual-rail lattice is $\Dfour=\{0,1\}^2$, ordered componentwise (four
values: bottom $(0,0)$, the two \emph{resolved} values $(1,0)$ and $(0,1)$
forming the face $R$, and top $(1,1)$). Define the involutions
$\sigma(a,b)=(1-a,\,1-b)$ (negation; order-reversing) and $P(a,b)=(b,a)$
(rail-swap; monotone). They commute and generate
$V_4=\{\mathrm{id},\sigma,P,\sigma P\}$.

\section{The pricing character}\label{sec:char}

$H^1(V_4;\F_2)=\operatorname{Hom}(V_4,\F_2)\cong\F_2^2$, with the three
nonzero characters determined by their kernels: $\chis$ (kernel
$\{\mathrm{id},P\}$), $\chip$ (kernel $\{\mathrm{id},\sigma\}$), and the
face-flip character $\chis+\chip$ (kernel $\{\mathrm{id},\sigma P\}$) ---
the character that is $1$ exactly on the elements moving the resolved face
$R$. For $\gamma\in V_4$ let $\nu(\gamma)\in\{0,1\}$ be its negation cost:
$0$ if $\gamma$ is monotone, $1$ otherwise.

\begin{proposition}[gate typing]\label{prop:gate}
As vectors over $(\mathrm{id},\sigma,P,\sigma P)$:
\begin{center}
\begin{tabular}{lcccc}
\toprule
$\gamma$ & monotone & $\nu(\gamma)$ & $\chis(\gamma)$ & $(\chis+\chip)(\gamma)$\\
\midrule
$\mathrm{id}$ & yes & $0$ & $0$ & $0$\\
$\sigma$      & no  & $1$ & $1$ & $1$\\
$P$           & yes & $0$ & $0$ & $\mathbf{1}$\\
$\sigma P$    & no  & $1$ & $1$ & $\mathbf{0}$\\
\bottomrule
\end{tabular}
\end{center}
Hence $\nu=\chis$ pointwise on $V_4$, and $\chis$ is the \textbf{unique}
nonzero character with this property: $\chip$ fails at $\sigma$, and the
face-flip character fails in both directions --- at $P$ (flip without cost)
and at $\sigma P$ (cost without flip).
\end{proposition}

\begin{proof}
Inspection of the four actions on $\Dfour$: $\sigma$ swaps both the
endpoints $(0,0)\leftrightarrow(1,1)$ and the face
$(1,0)\leftrightarrow(0,1)$ and reverses the order; $P$ fixes the endpoints,
swaps the face, and is monotone; $\sigma P$ swaps the endpoints, fixes the
face, and is not monotone (it exchanges bottom and top). Machine-checked as
\texttt{gate\_table} and \texttt{chiSigma\_eq\_not\_kMono} (axiom-free
kernel \texttt{decide}).
\end{proof}

The two witnesses deserve names: $P$ is the \emph{first separating witness}
(face-flip $1$, cost $0$) and $\sigma P$ the \emph{second} (cost $1$,
face-flip $0$). Together they show that ``pricing negation'' and ``moving
the resolved face'' are different characters --- they differ exactly by
$\chip$.

\section{Subadditivity and the lower bound}\label{sec:subadd}

\begin{lemma}[alternation]\label{lem:alt}
For every binary word $w$,
$\mathrm{rises}(w)\le\mathrm{drops}(w)+1$; if $w$ begins with $1$, then
$\mathrm{rises}(w)\le\mathrm{drops}(w)$.
\end{lemma}

\begin{proof}
Induction on the word, tracking the first letter: each maximal $1$-block
contributes at most one rise on its left and one drop on its right, and only
the initial $1$-block can lack a preceding rise. (Lean:
\texttt{rises\_le\_drops\_succ}.)
\end{proof}

\begin{lemma}[drop transfer]\label{lem:transfer}
Let $g$ be monotone in both arguments, $h$ arbitrary, and
$f(x)=g(x,\neg h(x))$. If $x<x'$ is a chain step with $f(x)=1$ and
$f(x')=0$, then $h(x)=0$ and $h(x')=1$ (an ascent of $h$).
\end{lemma}

\begin{proof}
If $\neg h(x)\le\neg h(x')$ then $(x,\neg h(x))\le(x',\neg h(x'))$ and
monotonicity of $g$ forces $f(x)\le f(x')$, a contradiction. So $\neg h$
descends, i.e.\ $h$ ascends.
\end{proof}

\begin{theorem}[subadditivity]\label{thm:subadd}
For monotone $g$ and arbitrary $h$,
$d\bigl(g(x,\neg h(x))\bigr)\le d(h)+1$. In fact, along any single chain,
$\mathrm{drops}(f\text{-word})\le\mathrm{rises}(h\text{-word})
\le\mathrm{drops}(h\text{-word})+1$; the argument uses only adjacent steps,
so it is valid over any relation, in particular any finite poset.
(Lean: \texttt{subadd\_chain}, \texttt{subadd\_decOn}.)
\end{theorem}

\begin{proof}
Each descent of $f$ is matched by an ascent of $h$ at the same step
(Lemma~\ref{lem:transfer}); ascents are bounded by descents plus one
(Lemma~\ref{lem:alt}); descents of $h$ along any chain are at most $d(h)$.
\end{proof}

\begin{corollary}[tower lower bound]\label{cor:lower}
Define $T_0=$ the monotone functions and
$T_{k+1}=\{\,g(x,\neg h(x)) : g\text{ monotone},\ h\in T_k\,\}$. Then
$f\in T_k$ implies $d(f)\le k$. (Lean: \texttt{F1\_lower}, via the
end-value-refined form
$\mathrm{rises}(w)\le\mathrm{drops}(w)+[w\text{ ends in }1]$.)
\end{corollary}

\begin{corollary}[no skipping]\label{cor:noskip}
If $f=g(x,\neg h)$ with $g$ monotone, then $d(h)\ge d(f)-1$.
\end{corollary}

\section{The three-zone witness}\label{sec:witness}

Fix $f$ with $d(f)=k\ge 1$. Call a pair $x<x'$ \emph{violated} if $f(x)=1$
and $f(x')=0$; write $V^-$ and $V^+$ for the lower and upper points of
violated pairs.

\begin{lemma}[potential jump]\label{lem:jump}
Across a violated pair, $\omega(x)+1\le\omega(x')$. In particular
$\omega(x)\le k-1$ and $\omega(x')\ge 1$. (Lean: \texttt{omega\_viol}.)
\end{lemma}

\begin{proof}
Take a cover-path from $x$ up to $x'$ and induct along it. If the first
step already descends ($f=1$ at its bottom, $f=0$ at its top), strictness of
$\omega$ across a descent gives the jump, and monotonicity of $\omega$
carries it to $x'$. Otherwise $f$ is still $1$ at the intermediate point and
the induction hypothesis applies from there, with monotonicity of $\omega$
across the first step.
\end{proof}

\begin{theorem}[three-zone witness]\label{thm:witness}
Define
\[
h\;:=\;1 \text{ on } \{\omega\ge k\},\qquad
\neg f \text{ on } \{1\le\omega\le k-1\},\qquad
0 \text{ on } \{\omega=0\}.
\]
Then \textup{(a)} $h$ covers every violated pair: $h=0$ on $V^-$ and $h=1$
on $V^+$; and \textup{(b)} $d(h)=k-1$ exactly. (Lean:
\texttt{hgen\_covering}, \texttt{hgen\_drops\_le}, \texttt{omega\_hgen\_le},
over an abstract descent potential; \texttt{cube\_descPotential} discharges
the potential axioms for the concrete $\omega$.)
\end{theorem}

\begin{proof}
(a) For $x\in V^-$: $\omega(x)\le k-1$ by Lemma~\ref{lem:jump}, so $h(x)$ is
either $0$ (zone $0$) or $\neg f(x)=0$ (middle zone). For $x'\in V^+$:
$\omega(x')\ge 1$, so $h(x')$ is either $\neg f(x')=1$ (middle) or $1$ (top
zone).

(b) Along any chain, $\omega$ is nondecreasing, so the chain traverses the
zones in the order $\{\omega=0\}\to\{1\le\omega\le k-1\}\to\{\omega\ge k\}$,
each a contiguous (possibly empty) stretch. $h$ is constant $0$, then
$\neg f$, then constant $1$. A descent of $h$ cannot start at the $0$-zone
($h=0$), cannot occur inside the top zone ($h=1$), and cannot occur at
either zone boundary (entering the middle from $0$ cannot descend; leaving
the middle to the top ends at $1$). So every descent of $h$ lies inside the
middle stretch and is an \emph{ascent} of $f$ there. Inside the middle
stretch, each descent of $f$ strictly increases $\omega$, and $\omega$ stays
within $[1,k-1]$; hence the stretch contains at most $k-2$ descents of $f$,
and by alternation at most $k-1$ ascents. So $d(h)\le k-1$; equality follows
from Corollary~\ref{cor:noskip} applied to the monotone-context
factorization produced in \S\ref{sec:tower}.
\end{proof}

\begin{remark}[the naive witness fails]\label{rem:naive}
The tempting witness $h_1:=\neg f\wedge[\omega\ge 1]$ is wrong already at
$k=1$. On three variables take $f(010)=f(111)=1$ and $f=0$ elsewhere: then
$d(f)=1$, but along the chain $000<100<110<111$ the potential jumps to $1$
at $110$ through the \emph{side} chain $000<010<110$, so $h_1$ takes values
$0,0,1,0$ --- one descent, though $k-1=0$. The top-zone cap ($h\equiv 1$ on
$\{\omega\ge k\}$) is exactly what repairs this. The example generalizes:
side-chains can raise $\omega$ without any descent on the chain being
traversed.
\end{remark}

\section{The nested tower theorem}\label{sec:tower}

\begin{lemma}[monotone extension]\label{lem:ext}
Let $h$ cover every violated pair of $f$. Then there is a monotone $g$ with
$f(x)=g(x,\neg h(x))$: define
\[
g(y,v)\;:=\;1 \iff \exists\,x\le y\ \bigl(h(x)=1\text{ or }v=1\bigr)
\wedge f(x)=1,
\]
a monotone function of $(y,v)$; the covering hypothesis rules out the only
conflict (a witness $x\le y$ with $f(x)=1$ forcing $g=1$ while $f(y)=0$ ---
such a pair is violated, so $h(x)=0$ and $h(y)=1$, and the disjunct
$h(x)\vee\neg h(y)$ fails). No choice principle is needed --- the extension
is an explicit finite disjunction. (Lean: \texttt{gExt}, \texttt{gExt\_mono},
\texttt{gExt\_agrees}.)
\end{lemma}

\begin{theorem}[nested tower]\label{thm:tower}
For every nonconstant $f:R^n\to R$,
\[
f\in T_k \iff d(f)\le k.
\]
Hence $d(f)$ is the minimal length of a nested $\chis$-tower over the
monotone clone, each rung consuming exactly one $\chis$-twist. (Lean:
\texttt{tower\_of\_decPts} for the upper bound --- constructive, producing
the witness of \S\ref{sec:witness} at each rung --- and \texttt{F1\_lower}
for the lower bound.)
\end{theorem}

\begin{proof}
($\Leftarrow$, upper bound) Induction on $k$. If $d(f)=0$, $f$ is monotone.
Otherwise let $h$ be the three-zone witness for $K:=d(f)$: by
Theorem~\ref{thm:witness}(b) $d(h)=K-1$, by induction $h\in T_{K-1}$, and by
Lemma~\ref{lem:ext} with Theorem~\ref{thm:witness}(a) $f=g(x,\neg h)$ for a
monotone $g$, so $f\in T_K\subseteq T_k$. ($\Rightarrow$ is
Corollary~\ref{cor:lower}.)
\end{proof}

\begin{remark}[models]\label{rem:models}
In the formula model this recovers $I_{\mathrm{for}}(f)=d(f)$
\cite{morizumi} with a \emph{nested} optimal form: Morizumi's construction
is the parallel decomposition $f_1\vee(f_2\wedge\neg f_t)$, in which the
recursive part and the fresh negation sit in separate branches;
Theorem~\ref{thm:tower} shows a single nested branch suffices. In the
circuit model the tower reading diverges from $I(f)$ once $d\ge 3$, by
Markov's logarithmic law --- the tower is a formula-model object.
\end{remark}

\begin{remark}[level, not degree]\label{rem:level}
The tower index is a filtration level at fixed cohomological degree one. No
identification with any higher-degree or higher-level operator theory is
claimed.
\end{remark}

\section{Mechanization}\label{sec:lean}

All statements are verified in core Lean~4 (no mathlib), 19 theorems; the
axiom print for every theorem is \texttt{[propext, Quot.sound]}, and the
finite character tables (\texttt{gate\_table},
\texttt{chiSigma\_eq\_not\_kMono}) depend on \textbf{no axioms}. There is no
\texttt{native\_decide} and no \texttt{sorry}; the upper bound is
constructive (no choice axiom). Correspondence:

\begin{center}
\begin{tabular}{ll}
\toprule
Paper & Lean (\texttt{lean/DecBridge/})\\
\midrule
Lemma~\ref{lem:alt} & \texttt{Words: rises\_le\_drops\_succ}\\
Lemma~\ref{lem:transfer} & \texttt{Subadd: drop\_transfer}\\
Theorem~\ref{thm:subadd} & \texttt{Subadd: subadd\_chain, subadd\_decOn}\\
Proposition~\ref{prop:gate} & \texttt{Gates: gate\_table, chiSigma\_eq\_not\_kMono}\\
Theorem~\ref{thm:witness} (abstract) & \texttt{Hgen: hgen\_covering, hgen\_drops\_le}\\
$\omega$ is a descent potential & \texttt{Cube: cube\_descPotential}\\
Lemma~\ref{lem:jump} & \texttt{CubeViol: omega\_viol}\\
$\omega=$ chain-max (both directions) & \texttt{CubeDec: chainDrops\_le\_omega, omega\_attained}\\
Theorem~\ref{thm:witness}(b) exact & \texttt{CubeDec: omega\_hgen\_le}\\
Lemma~\ref{lem:ext} & \texttt{Tower: gExt, gExt\_mono, gExt\_agrees}\\
Theorem~\ref{thm:tower} upper & \texttt{Tower: tower\_of\_decPts}\\
Cor.~\ref{cor:lower} / Thm.~\ref{thm:tower} lower & \texttt{Tower: F1\_lower, decPts\_le\_of\_tower}\\
\bottomrule
\end{tabular}
\end{center}

The Lean development is slightly \emph{more} general than the text:
subadditivity is proved for chains of an arbitrary relation (no order
axioms), and the witness theorem over an abstract descent potential, of
which the cube $\omega$ (defined by weight-fuel recursion over lower covers)
is an instance. The upper bound is stated over an explicit finite point
family (agreement on the listed points), which is the meaningful content on
the cube. Independent stdlib-only Python verifications (exhaustive $n\le 4$;
randomized $n=5$) accompany the artifact.

\subsection*{Authorship and provenance}

Produced by \emph{KoreoLoop}, an autonomous multi-agent research loop
developed and operated by the author, running frontier language models
(Anthropic Claude family) for discovery, formalization, and adversarial
verification, with the \textbf{Lean~4 kernel as the final acceptance gate}.
The author directed the research programme and verified the pipeline. In
line with the author's research-ethics policy, no claim of academic priority
is made beyond full disclosure of how the work was produced. Corrections are
invited. Artifact:
\url{https://github.com/ycmath/cohomological-price-of-not}.

\begin{thebibliography}{10}

\bibitem{markov} A.~A. Markov, On the inversion complexity of a system of
functions, \emph{J.~ACM} 5(4), 331--334, 1958. (Russian original:
\emph{Doklady AN SSSR} 116(6), 917--919, 1957.)

\bibitem{morizumi} H.~Morizumi, Limiting negations in formulas,
\emph{ICALP 2009}, LNCS 5555, 701--712. (Preliminary version:
arXiv:0811.0699, 2008.)

\bibitem{fischer} M.~J. Fischer, Lectures on network complexity, Tech.
Report 1104, Yale University, 1974 (revised 1996).

\bibitem{sw} M.~Santha, C.~Wilson, Limiting negations in constant depth
circuits, \emph{SIAM J.~Comput.} 22(2), 294--302, 1993.

\bibitem{st} S.~Sung, K.~Tanaka, Limiting negations in bounded-depth
circuits: an extension of Markov's theorem, \emph{ISAAC 2003}, LNCS 2906,
108--116.

\bibitem{am} K.~Amano, A.~Maruoka, A superpolynomial lower bound for a
circuit computing the clique function with at most
$\frac{1}{6}\log\log n$ negation gates, \emph{SIAM J.~Comput.} 35(1),
201--216, 2005.

\bibitem{km} V.~V. Kochergin, A.~V. Mikhailovich, Inversion complexity of
functions of multi-valued logic, arXiv:1510.05942, 2015; journal version:
\emph{Discrete Math.\ Appl.}, 2017.

\bibitem{gk} S.~Guo, I.~Komargodski, Negation-limited formulas,
\emph{APPROX/RANDOM 2015} (LIPIcs).

\bibitem{wcy1} W.~C. Yang, Finite-Energy Epistemic Logic with
Conservative Pointed Extension and Negation Geometry, public edition v1.0,
2026. \url{https://github.com/ycmath/finite-energy-epistemic-logic}
(an earlier draft circulated as ``Finite-energy epistemic logic with
T0-preserving open updates'').

\bibitem{wcy2} W.~C. Yang, The Price of NOT on D4, public edition v1.0,
2026. \url{https://github.com/ycmath/price-of-not-on-d4}

\end{thebibliography}

\end{document}
